\documentclass[11pt]{article}
\usepackage[margin=1in]{geometry}
\usepackage{amsmath,amssymb}
\usepackage{bm}
\usepackage{booktabs}
\usepackage{xcolor}
\usepackage{hyperref}
\hypersetup{colorlinks=true,linkcolor=blue,urlcolor=blue}

\newcommand{\rhat}{\hat{\mathbf r}}
\newcommand{\that}{\hat{\boldsymbol\theta}}
\newcommand{\tautt}{\texttt{earth\%tau}}

\title{Sensitivity of \textsc{global1d} solutions to the thin-sheet\\
surface conductance \tautt}
\author{Anna Kelbert with Claude}
\date{August 2026}

\begin{document}
\maketitle

\begin{abstract}
\textsc{global1d} (\texttt{field1d.f90}/\texttt{field1d\_s2.f90}) carries a single scalar parameter,
\tautt\ (units Siemens), representing an optional thin conducting sheet superimposed at the model's
outer surface $r=r_0$; we denote it $\tau$ below. This note derives how $\tau$ enters the governing
equations, identifies the natural dimensionless control parameter, and reports a numerical sensitivity
test against the closed-form analytic solution for a uniformly conducting sphere. We find (i) no
numerical advantage to a nonzero $\tau$: the solver is exact and well-behaved at $\tau=0$, which is
also the case the closed-form reference itself assumes; (ii) the error introduced by a nonzero $\tau$
scales linearly with the dimensionless parameter $Q=\omega\mu_0\tau\,r_0$ over eight orders of
magnitude, with the linear-regime coefficient implying the perturbation only reaches order unity near
$Q\sim40$; and (iii) a dedicated Earth-scale test (matching the conductor's own electrical size to the
toy sphere's, to isolate the effect of $Q$ alone) confirms this threshold and identifies the
mechanism beyond it: once $Q$ dominates the conductor's own pre-jump value, the boundary condition
is controlled by the artificial thin sheet rather than the layered conductor, so the response
\emph{saturates} in magnitude and can \emph{overshoot and reverse sign} in phase, rather than growing
without bound. (iv) A further check using the actual layered Earth conductivity model
(\texttt{LWS/layered\_GDE\_rho.prm}) at the production period of 3600\,s shows that, at the same $Q$,
the real Earth's much higher conductivity suppresses the perturbation well below what the
matched-regime toy-sphere comparison of (iii) predicted --- the $Q\sim40$ threshold is specific to
that toy conductor's own electrical regime, not a universal constant.
\end{abstract}

\section{How $\tau$ enters the equations}

Both solvers solve for the toroidal potential $T_l(r)$ via a transfer-matrix recursion through the
layered conductivity structure, matched to an insulating atmosphere above $r_0$. Ordinary continuity
of the layered-earth solution requires both $T$ and $T'$ continuous across every interface. At the
model's outer surface, \emph{both} solvers instead apply a modified matching condition
(field1d\_s2.f90 eq.~17; field1d.f90, algebraically identical, independently written):
\begin{equation}
  T_l'(r_0^+) \;=\; T_l'(r_0^-) \;-\; i\,\omega\,\mu_0\,\tau\,T_l(r_0),
  \qquad T_l(r_0^+) = T_l(r_0^-).
  \label{eq:jump}
\end{equation}
$T$ itself remains continuous; only $T'$ jumps, by a term proportional to $\tau$.

\paragraph{Physical interpretation.} This is the classical thin-sheet (Price 1949) approximation. A
sheet of conductance $\tau=\int\sigma\,dz$ carries an induced surface current $K=\tau\,E_{\rm
tangential}$ (Ohm's law integrated through the sheet's negligible thickness). Since $H_\theta,H_\phi
\propto T'/r$ and $E_\theta,E_\phi\propto T$, Amp\`ere's law relates a surface current directly to a
jump in tangential $H$ --- exactly \eqref{eq:jump}. Setting $\tau=0$ removes the sheet entirely and
recovers ordinary continuity.

\paragraph{Dimensionless control parameter.} Comparing the added term in \eqref{eq:jump} against the
natural scale of $T'$ itself ($T'/T\sim O(l)/r_0$) gives the classical thin-sheet parameter
\begin{equation}
  Q \;=\; \omega\,\mu_0\,\tau\,r_0 ,
  \label{eq:Q}
\end{equation}
which governs the size of the perturbation independent of the specific $(r_0,\omega)$ used in any
particular test or production run.

\section{Numerical test}

We use the same single homogeneous-sphere configuration as the existing closed-form regression test
(\texttt{test\_unit\_sphere.f90}: $r_0=1000$\,m, $\sigma=1$\,S/m, $\omega=1$\,rad/s, unit external
$(l{=}2,m{=}{+}1)$ multipole source), for which S2 already matches the analytic solution
\emph{exactly} at $\tau=0$ (ratio $=1+0i$ for all five field components). The $\tau=0$ output
therefore serves directly as the analytic reference; no closed-form values need to be re-derived here.
$H_r,H_\theta,H_\phi,E_\theta,E_\phi$ are extracted at one fixed grid point in the vacuum region
immediately above the sphere for every $\tau$.

\subsection{Magnitude error vs.\ $\tau$}

Table~\ref{tab:magerr} sweeps $\tau$ from 1\,S down to 0, reporting $Q$ (eq.~\ref{eq:Q}) and the
maximum relative error across all five field components against the $\tau=0$ reference.

\begin{table}[h]
\centering
\begin{tabular}{rrr}
\toprule
$\tau$ [S] & $Q=\omega\mu_0\tau r_0$ & max relative error (5 components) \\
\midrule
$1\times 10^{0}$  & $1.257\times 10^{-3}$ & $3.058\times 10^{-5}$ \\
$1\times 10^{-1}$ & $1.257\times 10^{-4}$ & $3.058\times 10^{-6}$ \\
$1\times 10^{-2}$ & $1.257\times 10^{-5}$ & $3.058\times 10^{-7}$ \\
$1\times 10^{-3}$ & $1.257\times 10^{-6}$ & $3.058\times 10^{-8}$ \\
$1\times 10^{-4}$ & $1.257\times 10^{-7}$ & $3.058\times 10^{-9}$ \\
$1\times 10^{-5}$ & $1.257\times 10^{-8}$ & $3.058\times 10^{-10}$ \\
$1\times 10^{-6}$ & $1.257\times 10^{-9}$ & $3.058\times 10^{-11}$ \\
$1\times 10^{-7}$ & $1.257\times 10^{-10}$ & $3.058\times 10^{-12}$ \\
$1\times 10^{-8}$ & $1.257\times 10^{-11}$ & $3.058\times 10^{-13}$ (roundoff floor) \\
$0$               & $0$                    & $0$ (exact, by construction) \\
\bottomrule
\end{tabular}
\caption{Magnitude error vs.\ $\tau$, referenced to the $\tau=0$ analytic solution, for the toy
sphere ($r_0=1000$\,m, $\sigma=1$\,S/m, $\omega=1$\,rad/s). Error scales linearly with $\tau$ (and
hence with $Q$) across eight decades, until double-precision roundoff floors out below
$\tau\sim10^{-6}$. The fit is $\text{error}\approx 0.024\,Q$.}
\label{tab:magerr}
\end{table}

\subsection{Magnitude ratio and phase shift of $H$ at specific $\tau$ values}

Tables~\ref{tab:phaseHr}--\ref{tab:phaseHph} report the magnitude ratio $|H(\tau)/H(0)|$ and phase
shift $\arg H(\tau)-\arg H(0)$ (degrees) of each magnetic field component in turn, for
$\tau=10^{-4},1,10^{2},2.5\times10^{3}$\,S.

\begin{table}[h]
\centering
\begin{tabular}{rrrr}
\toprule
$\tau$ [S] & $Q$ & $|H_r/H_{r,0}|$ & $\Delta\phi_{H_r}$ [$^\circ$] \\
\midrule
$10^{-4}$         & $1.257\times10^{-7}$ & 1.000000 & $1.7\times10^{-7}$ \\
$1$               & $1.257\times10^{-3}$ & 0.999998 & $1.75\times10^{-3}$ \\
$10^{2}$          & $0.1257$             & 0.999722 & $0.1743$ \\
$2.5\times10^{3}$ & $3.142$              & 0.964484 & $3.086$ \\
\bottomrule
\end{tabular}
\caption{$H_r$ magnitude ratio and phase shift relative to $\tau=0$, for the toy sphere
($r_0=1000$\,m, $\sigma=1$\,S/m, $\omega=1$\,rad/s).}
\label{tab:phaseHr}
\end{table}

\begin{table}[h]
\centering
\begin{tabular}{rrrr}
\toprule
$\tau$ [S] & $Q$ & $|H_\theta/H_{\theta,0}|$ & $\Delta\phi_{H_\theta}$ [$^\circ$] \\
\midrule
$10^{-4}$         & $1.257\times10^{-7}$ & 1.000000 & $-3\times10^{-8}$ \\
$1$               & $1.257\times10^{-3}$ & 1.000000 & $-3.00\times10^{-4}$ \\
$10^{2}$          & $0.1257$             & 1.000051 & $-0.02992$ \\
$2.5\times10^{3}$ & $3.142$              & 1.006426 & $-0.5060$ \\
\bottomrule
\end{tabular}
\caption{$H_\theta$ magnitude ratio and phase shift relative to $\tau=0$, for the toy sphere
($r_0=1000$\,m, $\sigma=1$\,S/m, $\omega=1$\,rad/s).}
\label{tab:phaseHth}
\end{table}

\begin{table}[h]
\centering
\begin{tabular}{rrrr}
\toprule
$\tau$ [S] & $Q$ & $|H_\phi/H_{\phi,0}|$ & $\Delta\phi_{H_\phi}$ [$^\circ$] \\
\midrule
$10^{-4}$         & $1.257\times10^{-7}$ & 1.000000 & $-3\times10^{-8}$ \\
$1$               & $1.257\times10^{-3}$ & 1.000000 & $-3.00\times10^{-4}$ \\
$10^{2}$          & $0.1257$             & 1.000051 & $-0.02992$ \\
$2.5\times10^{3}$ & $3.142$              & 1.006426 & $-0.5060$ \\
\bottomrule
\end{tabular}
\caption{$H_\phi$ magnitude ratio and phase shift relative to $\tau=0$, for the toy sphere
($r_0=1000$\,m, $\sigma=1$\,S/m, $\omega=1$\,rad/s).}
\label{tab:phaseHph}
\end{table}

Tables~\ref{tab:phaseHth} and~\ref{tab:phaseHph} are numerically identical. This is not a coincidence:
$H_\theta$ and $H_\phi$ both derive from the same radial factor $T'(r)$, whose complex ratio to the
$\tau=0$ case does not depend on angular position, so the two components (evaluated at the same
radius, different angles) necessarily agree. $H_r$ (Table~\ref{tab:phaseHr}), driven by $T(r)$ rather
than $T'(r)$, follows an independent (smaller-magnitude, opposite-sign) trend.

\section{Earth-scale test: is $Q$ alone sufficient, and why does the response become nonlinear?}
\label{sec:earth}

The toy sphere has $r_0=1000$\,m, $\sigma=1$\,S/m, $\omega=1$\,rad/s (as in \S2); the real Earth has
$r_0\approx6.371\times10^6$\,m ($6371\times$ larger). Since $Q\propto r_0$ (eq.~\ref{eq:Q}), the
\emph{same} $\tau$ produces a vastly larger $Q$ at Earth scale --- Table~\ref{tab:earthQ} evaluates
this at the NA2026 project's own periods (162\,s, 3600\,s).

\begin{table}[h]
\centering
\begin{tabular}{r rr}
\toprule
$\tau$ [S] & $Q_{\rm Earth}$, 162\,s & $Q_{\rm Earth}$, 3600\,s \\
\midrule
$10^{-4}$ & $3.1\times10^{-5}$ & $1.4\times10^{-6}$ \\
$1$       & $0.31$ & $0.014$ \\
$10^{2}$  & $31$ & $1.4$ \\
$2.5\times10^{3}$ & $776$ & $35$ \\
\bottomrule
\end{tabular}
\caption{$Q$ at Earth scale for the same $\tau$ values as Tables~\ref{tab:phaseHr}--\ref{tab:phaseHph}.}
\label{tab:earthQ}
\end{table}

\paragraph{Methodology: isolating $Q$ from the conductor's own regime.} Naively re-running the toy
sphere with $r_0=6.371\times10^6$\,m but the same $\sigma=1$\,S/m, $\omega=1$\,rad/s would \emph{also}
push the conductor's own electrical size $|k_l|r_0=\sqrt{\omega\mu_0\sigma}\,r_0$ from $\approx1.12$
(order-1 induction, the toy sphere's regime) to $\approx1400$ (extreme skin effect) --- a second,
unrelated regime change that would confound any observed nonlinearity: we would no longer know
whether a deviation from linear-in-$Q$ behavior is caused by $\tau$ growing large, or by the
conductor itself behaving completely differently. To isolate the $\tau$/$Q$ effect cleanly, $\sigma$
at each Earth-scale period is instead chosen so that $|k_l|r_0$ \emph{equals} the toy sphere's own
value ($1.121$), holding the conductor's regime fixed while only $r_0$ (and hence $Q$, for the same
$\tau$) changes. This gives $\sigma=6.35\times10^{-7}$\,S/m at 162\,s and $\sigma=1.41\times10^{-5}$\,S/m
at 3600\,s --- unrealistically resistive for an actual whole-Earth sphere, but that is expected and
intentional: this is a controlled numerical experiment isolating one mechanism, not a production-accurate
Earth model. $H_\theta=H_\phi$ again (same argument as \S2.2), so only one column pair is shown for them.

\begin{table}[h]
\centering
\begin{tabular}{rr rr rr}
\toprule
$\tau$ [S] & $Q$ & $|H_r/H_{r,0}|$ & $\Delta\phi_{H_r}$ [$^\circ$]
           & $|H_{\theta,\phi}/H_{\theta,\phi,0}|$ & $\Delta\phi_{H_{\theta,\phi}}$ [$^\circ$] \\
\midrule
$10^{-4}$         & $3.11\times10^{-5}$ & 1.000000 & $4.3\times10^{-5}$  & 1.000000 & $-7.4\times10^{-6}$ \\
$1$               & $0.3105$            & 0.999065 & $0.4254$            & 1.000173 & $-0.07295$ \\
$10^{2}$          & $31.05$             & 0.882199 & $0.9808$            & 1.020275 & $-0.13990$ \\
$2.5\times10^{3}$ & $776.3$             & 0.878969 & $-0.1977$           & 1.020790 & $0.03518$ \\
\bottomrule
\end{tabular}
\caption{$H$ magnitude ratio and phase shift relative to $\tau=0$, Earth radius, period 162\,s,
$\sigma=6.35\times10^{-7}$\,S/m (matched $|k_l|r_0$).}
\label{tab:earth162}
\end{table}

\begin{table}[h]
\centering
\begin{tabular}{rr rr rr}
\toprule
$\tau$ [S] & $Q$ & $|H_r/H_{r,0}|$ & $\Delta\phi_{H_r}$ [$^\circ$]
           & $|H_{\theta,\phi}/H_{\theta,\phi,0}|$ & $\Delta\phi_{H_{\theta,\phi}}$ [$^\circ$] \\
\midrule
$10^{-4}$         & $1.40\times10^{-6}$ & 1.000000 & $1.9\times10^{-6}$ & 1.000000 & $-3.3\times10^{-7}$ \\
$1$               & $0.01397$           & 0.999976 & $0.01932$          & 1.000004 & $-0.00332$ \\
$10^{2}$          & $1.397$             & 0.989809 & $1.7626$           & 1.001871 & $-0.29856$ \\
$2.5\times10^{3}$ & $34.93$             & 0.881538 & $0.8515$           & 1.020380 & $-0.12055$ \\
\bottomrule
\end{tabular}
\caption{$H$ magnitude ratio and phase shift relative to $\tau=0$, Earth radius, period 3600\,s,
$\sigma=1.41\times10^{-5}$\,S/m (matched $|k_l|r_0$).}
\label{tab:earth3600}
\end{table}

\paragraph{Refining the linear-regime threshold.} Table~\ref{tab:magerr}'s fit, $\text{error}\approx
0.024\,Q$, implies the perturbation only reaches order unity once $Q\sim1/0.024\approx42$ --- not
$Q\sim1$ as the crude order-of-magnitude argument in \S1 suggested. This is exactly consistent with
Tables~\ref{tab:earth162}--\ref{tab:earth3600}: at $Q=0.31$--$1.4$ the response is still closely
linear in $Q$ (e.g.\ $\Delta\phi_{H_r}$ at $Q=0.31$ (162\,s) and $Q=1.4$ (3600\,s) both fall
close to the toy-sphere fit scaled to that $Q$), while by $Q=31$--$35$ (comparable to the refined
$Q\sim42$ threshold) the magnitude ratio has already dropped by ${\sim}12\%$ --- clearly nonlinear.

\paragraph{Why the response saturates rather than diverging.} From eq.~\eqref{eq:jump}, with
$T_l(r_0)=1$ by construction, $T_l'(r_0^+)=t_{np}(l)-iQ$, where $t_{np}(l)$ is the conductor's own
($\tau$-independent) pre-jump value. Once $Q\gg|t_{np}(l)|$, this sum is dominated by the $-iQ$ term
regardless of the conductor's own properties: $T_l'(r_0^+)\to-iQ$, a fixed $-90^\circ$ phase relative
to $T_l(r_0)$, independent of $\tau$ once $\tau$ is large enough. This is why the magnitude ratio
\emph{saturates} (Tables~\ref{tab:earth162}--\ref{tab:earth3600}: essentially unchanged between
$Q=31$ and $Q=776$, and between $Q=1.4$ and $Q=35$) rather than continuing to grow, and why the phase
shift can \emph{overshoot and reverse sign} (162\,s: $+0.98^\circ$ at $Q=31$ to $-0.20^\circ$ at
$Q=776$; 3600\,s: $+1.76^\circ$ at $Q=1.4$ down to $+0.85^\circ$ at $Q=35$) rather than growing
monotonically: for very large $\tau$ the physics is no longer governed by the layered conductor at
all, but entirely by the artificial thin sheet, and further increasing $\tau$ has diminishing effect
on the already-$\tau$-dominated boundary condition. This is a genuine physical saturation, not a
numerical artifact: the same qualitative behavior (rise, peak/plateau, mild reversal) appears at both
Earth-scale periods, at $Q$ values that differ by a factor of ${\sim}20$, using two independently
matched conductor regimes.

\section{Realistic layered Earth conductivity model}
\label{sec:gde}

\S\ref{sec:earth}'s Earth-scale test used a homogeneous sphere with $\sigma$ chosen only to match
the toy sphere's own $|k_l|r_0$ --- a controlled experiment isolating the $Q$ mechanism, not a
realistic Earth. Here we repeat the same $\tau$ sweep at Earth radius and the NA2026 project's
3600\,s period, but with the \emph{actual} layered conductivity profile used in production runs,
\texttt{LWS/layered\_GDE\_rho.prm}, read the same way \texttt{FWD1D.f90} builds \texttt{earth\%layer}/
\texttt{earth\%sigma} from a \texttt{layers}-format model file: 32 layers given as depth below the
surface (km) and $\log_{10}\rho$ ($\rho$ in $\Omega\cdot$m), converted to $\sigma=10^{-\log_{10}\rho}$,
plus a hardcoded $\sigma=10^5$\,S/m ``core'' layer beneath the deepest defined depth (3500\,km).
Table~\ref{tab:gdemodel} quotes every depth/$\sigma$ pair used.

\begin{table}[h]
\centering
\begin{minipage}{0.48\textwidth}
\centering
\begin{tabular}{rrr}
\toprule
depth [km] & $\log_{10}\rho$ & $\sigma$ [S/m] \\
\midrule
10   & 2.000000  & $1.000\times10^{-2}$ \\
40   & 3.000000  & $1.000\times10^{-3}$ \\
110  & 3.000000  & $1.000\times10^{-3}$ \\
160  & 2.522879  & $3.000\times10^{-3}$ \\
210  & 2.154902  & $7.000\times10^{-3}$ \\
260  & 1.745627  & $1.796\times10^{-2}$ \\
310  & 1.667885  & $2.148\times10^{-2}$ \\
360  & 1.688480  & $2.049\times10^{-2}$ \\
410  & 1.709952  & $1.950\times10^{-2}$ \\
460  & 1.699400  & $1.998\times10^{-2}$ \\
510  & 1.645104  & $2.264\times10^{-2}$ \\
560  & 1.540801  & $2.879\times10^{-2}$ \\
610  & 1.379495  & $4.174\times10^{-2}$ \\
660  & 1.153394  & $7.024\times10^{-2}$ \\
710  & 0.869989  & $1.349\times10^{-1}$ \\
760  & 0.593626  & $2.549\times10^{-1}$ \\
\bottomrule
\end{tabular}
\end{minipage}
\hfill
\begin{minipage}{0.48\textwidth}
\centering
\begin{tabular}{rrr}
\toprule
depth [km] & $\log_{10}\rho$ & $\sigma$ [S/m] \\
\midrule
800  & 0.430060  & $3.715\times10^{-1}$ \\
850  & 0.279000  & $5.260\times10^{-1}$ \\
900  & 0.279000  & $5.260\times10^{-1}$ \\
1000 & $-0.228100$ & $1.691$ \\
1100 & $-0.228100$ & $1.691$ \\
1200 & $-0.228100$ & $1.691$ \\
1300 & $-0.228100$ & $1.691$ \\
1400 & $-0.228100$ & $1.691$ \\
1500 & $-0.228100$ & $1.691$ \\
1620 & $-0.228100$ & $1.691$ \\
1764 & $-0.228100$ & $1.691$ \\
1936.8 & $-0.228100$ & $1.691$ \\
2144.16 & $-0.228100$ & $1.691$ \\
2400 & $-0.228100$ & $1.691$ \\
2900 & $-1.000000$ & $10.00$ \\
3500 & $-2.000000$ & $100.0$ \\
(core, $r<r_0-3500$\,km) & --- & $1.000\times10^5$ \\
\bottomrule
\end{tabular}
\end{minipage}
\caption{Layer depths and conductivities of the layered Earth model, \texttt{LWS/layered\_GDE\_rho.prm}
(32 layers, $\sigma=10^{-\log_{10}\rho}$, plus the hardcoded core layer), used for
Table~\ref{tab:gderesults}.}
\label{tab:gdemodel}
\end{table}

\begin{table}[h]
\centering
\begin{tabular}{rr rr rr}
\toprule
$\tau$ [S] & $Q$ & $|H_r/H_{r,0}|$ & $\Delta\phi_{H_r}$ [$^\circ$]
           & $|H_{\theta,\phi}/H_{\theta,\phi,0}|$ & $\Delta\phi_{H_{\theta,\phi}}$ [$^\circ$] \\
\midrule
$10^{-4}$         & $1.40\times10^{-6}$ & 1.000000 & $6\times10^{-8}$   & 1.000000 & $-1\times10^{-8}$ \\
$1$               & $0.01397$           & 0.999987 & $6.37\times10^{-4}$ & 1.000002 & $-9.56\times10^{-5}$ \\
$10^{2}$          & $1.397$             & 0.998695 & $0.05833$          & 1.000201 & $-0.00872$ \\
$2.5\times10^{3}$ & $34.93$             & 0.979892 & $0.02557$          & 1.003066 & $-0.00138$ \\
\bottomrule
\end{tabular}
\caption{$H$ magnitude ratio and phase shift relative to $\tau=0$, Earth radius, period 3600\,s,
\emph{realistic} layered conductivity (Table~\ref{tab:gdemodel}) rather than the $|k_l|r_0$-matched
homogeneous sphere of Table~\ref{tab:earth3600}. Same $\tau,Q$ values as every other table in this
report.}
\label{tab:gderesults}
\end{table}

\paragraph{The perturbation is markedly smaller than the matched-regime toy comparison predicted.}
$Q$ depends only on $\omega$, $\tau$, $r_0$ (eq.~\ref{eq:Q}), not on conductivity, so
Table~\ref{tab:gderesults} uses \emph{exactly} the same $Q$ values as Table~\ref{tab:earth3600}. Yet
at $Q=34.9$ ($\tau=2500$\,S), the realistic model gives $|H_r/H_{r,0}|=0.980$ and
$\Delta\phi_{H_r}=0.026^\circ$, versus $0.882$ and $0.85^\circ$ for the matched-regime homogeneous
sphere --- roughly $6\times$ smaller in both magnitude deviation and phase shift at the identical
$Q$. The mechanism from the previous paragraph explains why: saturation sets in once
$Q\gg|t_{np}(l)|$, and $|t_{np}(l)|$ --- the conductor's own pre-jump value --- is set by the
conductivity structure, not by $Q$. The real Earth model is far more conductive through most of its
depth (up to $\sigma\approx1.7$\,S/m in the mantle, $\sigma=10^5$\,S/m in the core; see
Table~\ref{tab:gdemodel}) than the thin, $|k_l|r_0$-matched homogeneous sphere used in
\S\ref{sec:earth} to isolate the $Q$ effect, giving it a correspondingly larger $|t_{np}(l)|$ and
therefore requiring a larger $Q$ before the $-iQ$ term in the jump condition dominates. The phase
shift is still non-monotonic in $Q$ here too ($0.058^\circ$ at $Q=1.4$, falling to $0.026^\circ$ at
$Q=34.9$) --- the same qualitative saturation/overshoot behavior identified in \S\ref{sec:earth}, just
at a much smaller absolute scale and a lower onset $Q$ than that section's toy-conductor threshold.
\textbf{Conclusion: the $Q\sim40$ threshold derived in \S2.1/\S\ref{sec:earth} is specific to that toy
conductor's own electrical regime, not a universal property of $Q$ alone} --- a genuinely more
conductive Earth model tolerates a larger $\tau$ (at fixed $r_0,\omega$) before departing from linear
response. This makes the practical recommendation below (set $\tau=0$) even safer for production runs
than \S\ref{sec:earth} alone would suggest: at the current default $\tau=10^{-4}$\,S, $Q\approx1.4\times
10^{-6}$ at 3600\,s regardless of which conductivity model is used, many orders of magnitude below
where any version of this nonlinearity begins.

\section{Discussion}

\paragraph{No numerical advantage to nonzero $\tau$.} $\tau$ appears in exactly one place in each
solver: the additive term in \eqref{eq:jump}. It is never a divisor, matrix pivot, or regularizer, so
$\tau=0$ introduces no degeneracy --- consistent with the exact match already obtained at
$\tau=0$ in the existing regression test. There is no floating-point or conditioning reason to
prefer a small nonzero value over exactly zero.

\paragraph{Practical recommendation.} For runs that are not deliberately modeling a real conductive
surface sheet, set $\tau=0$: it is the mathematically clean choice, matches the validated exact-match
test case, and removes an otherwise-unexplained free parameter. \texttt{FWD1D.f90}'s current default
$\tau=10^{-4}$ has no numerical justification we could identify and can be replaced by $0$ without
any accuracy cost.

\paragraph{Scale dependence.} Since $Q\propto r_0$ (eq.~\ref{eq:Q}), the toy sphere's $Q$ values
($\tau=2500$\,S gives only $Q=3.14$) do \emph{not} directly represent Earth-scale behavior --- at
Earth's radius the same $\tau$ values give $Q$ up to $776$ (Table~\ref{tab:earthQ}), well beyond
where Table~\ref{tab:magerr}'s linear fit holds. \S\ref{sec:earth} evaluates this directly rather
than extrapolating: at $\tau=10^{-4}$\,S, $Q_{\rm Earth}\sim10^{-5}$--$10^{-6}$, confirming the
current default is negligible at Earth scale too and can safely be replaced by $0$. At
$\tau=100$--$2500$\,S, $Q_{\rm Earth}$ reaches the regime where \S\ref{sec:earth}'s toy-conductor
comparison first showed nonlinearity ($Q\gtrsim40$) --- but \S\ref{sec:gde} shows this specific
threshold is conductor-dependent, not a property of $Q$ alone: repeating the same $\tau$ sweep with
the actual production conductivity model (\texttt{LWS/layered\_GDE\_rho.prm}) at 3600\,s gives a
perturbation roughly $6\times$ smaller, at the identical $Q$, than the matched-regime toy sphere
predicted, because the real Earth's higher conductivity gives it a correspondingly larger
$|t_{np}(l)|$. This is not merely a mechanism demonstration but a production-representative estimate
already obtained (\S\ref{sec:gde}, Table~\ref{tab:gderesults}): if a nonzero $\tau$ representing a
genuine conductive surface feature (e.g.\ the ocean) is intentionally used at Earth scale, its
quantitative effect on this specific conductivity model is now known directly, not merely bounded by
the toy-sphere mechanism study.

\end{document}
